% !TEX root = MathPhysExercises.tex
\documentclass[12pt,a4paper]{article}
\usepackage{preamble}

\renewcommand{\notesuniversity}{Санкт-Петербургский государственный университет}
\renewcommand{\notesfaculty}{Факультет математики и механики}
\renewcommand{\notestitle}{Математическая физика}
\renewcommand{\notessubtitle}{Упражнения}
\renewcommand{\notesauthor}{}
\renewcommand{\notesdate}{}
\renewcommand{\notesheader}{Математическая физика \(\cdot\) Упражнения}

\hypersetup{pdftitle={Математическая физика — Упражнения}}

\theoremstyle{definition}
\newtheorem{notesexercise}{Упражнение}[section]
\makeatletter
\renewcommand{\exr}[2]{\notes@wrap{notesexercise}{#1}{#2}}
\makeatother

\begin{document}
\notesmaketitle

\tableofcontents

\newpage

% =====================================================================
\section{Практика 1}
\renewcommand{\notesheader}{Математическая физика \(\cdot\) Практика 1 \(\cdot\) Упражнения}
% =====================================================================

\exr{}{
Пусть \(X\), \(Y\) --- топологические векторные пространства,
\(X \hookrightarrow Y\), где \(\hookrightarrow\) такое, что
\(\exists A \colon X \to Y\), где \(A\) линейно, непрерывно и инъективно,
при этом \(X\) всюду плотно в \(Y\). Тогда \(Y^* \hookrightarrow X^*\).
}

\rem{Контекст}{
Сходимость в \(\DD(\Omega)\):
\(\phi_k \to \phi\) тогда и только тогда, когда
\begin{enumerate}[leftmargin=1.6em,itemsep=3pt]
    \item \(\forall \alpha\ \ D^\alpha \phi_k \rightrightarrows D^\alpha \phi\);
    \item \(\exists K\) --- компакт: \(\forall k\ \supp \phi_k \subset K \subset \Omega\).
\end{enumerate}
}

\exr{}{
Привести 2 топологических пространства с одинаковой сходимостью
(Богачев В.\,И., ``Основы теории мер'', Том~2, стр.~70).
}

\exr{}{
Показать, что такая сходимость не может быть порождена метрикой.
}

\rem{Контекст}{
Локально конечный борелевский заряд \(\mu \in \mathcal{M}_{\loc}(\Omega)\)
сопоставляют функционалу
\[
\langle \widetilde{\mu},\, \phi \rangle
:=
\int_\Omega \phi\,d\mu.
\]
}

\exr{}{
\(\widetilde{\mu} \in \DD'(\Omega)\) и
\(\mathcal{M}_{\loc}(\Omega) \to \DD'(\Omega):\ \mu \mapsto \widetilde{\mu}\).
\[
\langle \mu,\, \phi \rangle
=
\int_\Omega \phi\,d\mu.
\]
}

\rem{Контекст}{
Регулярная обобщённая функция --- это функционал \(\widetilde{f}\),
отвечающий \(f \in L_1^{\loc}(\Omega)\):
\[
\langle \widetilde{f},\, \phi \rangle
=
\int_\Omega f\,\phi.
\]
Дельта-функция задаётся формулой \(\langle \delta,\, \phi \rangle = \phi(0)\).
}

\exr{}{\(\delta\) --- не регулярная функция.}

\rem{Контекст}{
\(\supp f\) --- наименьшее замкнутое множество, вне которого \(f = 0\).
Равенство \(f = 0\) на открытом \(\Omega' \subset \Omega\) означает, что
\(\langle f,\, \phi \rangle = 0\) для всех \(\phi \in \DD(\Omega)\)
с \(\supp\phi \subset \Omega'\).
}

\exr{}{
Если \(f\) --- заряд или регулярная функция, то \(\supp f\) в этом смысле
совпадает с \(\supp f\) в стандартном смысле.
}

\exr{}{
\(\eta \in C^\infty(\Omega)\) и \(\eta \equiv 1\) на \(\supp f\),
\(f \in \DD(\Omega)\). Тогда \(\eta f = f\).
}

\rem{Контекст}{
Замена переменных вдоль диффеоморфизма \(\Phi \colon \Omega \to \Omega_1\):
\[
\langle f \circ \Phi,\, \phi \rangle
:=
\biggl\langle
    f,\,
    \frac{\phi \circ \Phi^{-1}}{|\det \Phi'|}
\biggr\rangle.
\]
Умножение на \(a \in C^\infty\):
\(\langle af,\, \phi \rangle = \langle f,\, a\phi \rangle\).
}

\exr{}{
\(a \in C^\infty(\Omega)\), \(f \in \DD'(\Omega)\),
\(\Phi\) --- диффеоморфизм \(\Omega \to \Omega_1\). Тогда
\[
(af)\bigl(\phi(x)\bigr)
=
a\bigl(\phi(x)\bigr)\,f\bigl(\phi(x)\bigr).
\]
}

\rem{Контекст}{
Обобщённая производная:
\[
\langle D^\alpha f,\, \phi \rangle
=
(-1)^{|\alpha|}\,
\langle f,\, D^\alpha\phi \rangle.
\]
}

\exr{}{
\(D^\alpha D^\beta f = D^\beta D^\alpha f = D^{\alpha+\beta} f\).
}

\exr{}{
Правило Лейбница
\(D_i(a \cdot f) = D_i a \cdot f + a D_i f\).
}

\rem{Контекст}{
\(\operatorname{sing\,supp} f\) --- наименьшее замкнутое множество,
вне которого \(f\) совпадает с какой-то \(C^\infty\)-функцией.
Носитель \(\supp f\) --- как выше: наименьшее замкнутое множество,
вне которого \(f = 0\).
}

\exr{}{
\(\supp D^\alpha f \subset \supp f\) и привести пример, когда нет равенства.
}

\exr{}{
\(\operatorname{sing\,supp} D^\alpha f \subset \operatorname{sing\,supp} f\)
и привести пример, когда нет равенства.
}

\newpage
% =====================================================================
\section{Практика 2}
\renewcommand{\notesheader}{Математическая физика \(\cdot\) Практика 2 \(\cdot\) Упражнения}
% =====================================================================

\rem{Контекст}{
Функция Хевисайда \(\Theta(x) = \chi_{[0,+\infty)}(x)\)
удовлетворяет \(\Theta' = \delta\) в \(\DD'(\R)\).
}

\exr{Кусочно-гладкая функция}{
Пусть \(f\) --- кусочно-гладкая функция на \(\R\).
Найти \(f' \in \DD'(\R)\).
\par\smallskip
\noindent\textit{Подсказка:} понадобится функция Хевисайда.
\par\medskip
\centering
\begin{tikzpicture}[x=0.72cm, y=0.52cm, baseline=(current bounding box.center)]
  \draw[-{Stealth}, gray!70] (-2.6,0) -- (3.4,0);
  \draw[-{Stealth}, gray!70] (0,-0.35) -- (0,2.1);
  \draw[black, thick] (-2.3,0.35) -- (-0.15,1.05);
  \draw[black, thick] (0.15,1.55) -- (3.1,0.55);
  \fill[black] (-0.15,1.05) circle (1.6pt);
  \fill[black] (0.15,1.55) circle (1.6pt);
  \draw[black, densely dashed] (-0.15,1.05) -- (0.15,1.05);
\end{tikzpicture}\par
}

\rem{Контекст}{
\(\widetilde{f}\) --- регулярная обобщённая функция, отвечающая
\(f \in L_1^{\loc}\). Сдвиг понимается как замена переменных,
производная --- как выше. В частности,
\[
\delta'(x)
=
\lim_{h \to 0}
\frac{\delta(x+h) - \delta(x)}{h}.
\]
}

\exr{Разностная производная}{
Показать, что
\[
\widetilde{f}'_{x_i}
=
\lim_{h \to 0}
\frac{\widetilde{f}(x + h e_i) - \widetilde{f}(x)}{h}.
\]
}

\rem{Контекст}{
Коммутативное ассоциативное произведение обобщённых функций,
согласованное с умножением на \(C^\infty\), невозможно.
Поскольку \(x\,\delta(x) = 0\),
\[
\bigl(\delta(x) \cdot x\bigr)\, \Pvf\frac{1}{x} = 0,
\]
а из тождества \(x\, \Pvf\frac{1}{x} = 1\) следует
\[
\delta(x) \cdot \Bigl(x\, \Pvf\frac{1}{x}\Bigr) = \delta(x).
\]
}

\exr{Дизъюнктные носители}{
Если \(\supp f \cap \supp g = \varnothing\), то такое произведение
определить можно.
}

\newpage
% =====================================================================
\section{Практика 3}
\renewcommand{\notesheader}{Математическая физика \(\cdot\) Практика 3 \(\cdot\) Упражнения}
% =====================================================================

\rem{Контекст}{
Прямое произведение \(f(x)g(y)\) коммутативно на разложимых
пробных функциях \(\phi(x)\psi(y)\), а значит и на их линейных
комбинациях. Дальше доказывается, что такие линейные комбинации
плотны в \(\DD(\Omega_1 \times \Omega_2)\). В конспекте это
рассуждение записано при \(\dim \Omega_1 = \dim \Omega_2 = 1\).
}

\exr{}{Доказать для общего случая.}

\rem{Контекст}{
Сходимость в \(\DD'(\Omega)\):
\(f_k \to f\), если
\(\langle f_k,\, \phi \rangle \to \langle f,\, \phi \rangle\)
для всякой \(\phi \in \DD(\Omega)\).
Пределы ниже понимаются в этой сходимости.
}

\exr{}{
\(\omega_\epsilon\) --- аппроксимативная единица.
\(\displaystyle\lim_{\epsilon \to 0} \omega_\epsilon\)
(ответ: \(\delta\)-функция).
}

\exr{}{
\(\displaystyle\lim_{h \to 0} \frac{f(x+h) - f(x)}{h} = ?\)
(дома).
}

\exr{}{
\(\displaystyle\lim_{\epsilon \to 0} \frac{1}{x + i\epsilon} = ?\)
в \(\DD'(\R)\)
(решили на паре).
}

\rem{Контекст}{
В \(\DD'(\R)\)
\[
\frac{1}{x + i\epsilon}
\longrightarrow
\Pvf\frac{1}{x} - i\pi\delta.
\]
}

\exr{}{Доказать это для \(\dfrac{1}{x - i\epsilon}\).}

\rem{Контекст}{
\(\Pvf\) --- главное значение: для сингулярного ядра
\[
\Bigl\langle
    \Pvf \frac{A\bigl(\tfrac{x}{|x|}\bigr)}{|x|^n},\,
    \phi
\Bigr\rangle
=
\vp \int_{\R^n}
    \frac{A\bigl(\tfrac{x}{|x|}\bigr)}{|x|^n}\,
    \phi(x)\,dx,
\]
где интеграл в смысле главного значения --- предел интегралов
по дополнению шара \(B_\epsilon\). Пределы ниже --- в \(\DD'(\R)\).
}

\exr{}{
\[
\lim_{n \to \infty}
\Pvf\frac{\cos nx}{x}
= ?,
\]
\[
\lim_{n \to \infty}
\Pvf\frac{\sin nx}{x}
= ?.
\]
}

\newpage
% =====================================================================
\section{Практика 4}
\renewcommand{\notesheader}{Математическая физика \(\cdot\) Практика 4 \(\cdot\) Упражнения}
% =====================================================================

\rem{Контекст}{
Свёртка \(f, g \in \DD'(\R^n)\), если предел существует и не зависит
от финитизатора \(\{\eta_k\}\) (\(\eta_k \to 1\): на каждом компакте
\(\eta_k\) с какого-то момента тождественно равна \(1\), и
\(|D^\alpha \eta_k| \le C_\alpha\)):
\[
\langle f * g,\, \phi \rangle
:=
\lim_k
\langle f(x)\, g(y),\, \eta_k(x,y)\, \phi(x+y) \rangle.
\]
Сдвиг обобщённой функции --- замена переменных \(x \mapsto x + h\).
}

\rem{Контекст}{
Если свёртка существует, то
\[
D^\alpha(f * g) = (D^\alpha f) * g = f * (D^\alpha g),
\]
и эти свёртки существуют и равны.
}

\exr{}{
Привести пример, когда \(\exists\, D_i f * g\), \(\exists\, f * D_i g\),
но они не равны.
}

\exr{}{
\(f(x) * g(x+h) = (f * g)(x+h)\) --- оператор свёртки коммутирует
с оператором сдвига.
}

\rem{Контекст}{
Доказано включение
\(\supp(f * g) \subset \overline{\supp f + \supp g}\).
Без замыкания сумма \(\supp f + \supp g\) не обязательно замкнута.
}

\exr{}{
\(\bigl(\overline{\supp f + \supp g}\bigr)^c\) не обязательно открыто.
}

\exr{}{
Пример \(f\) и \(g\):
\(\supp(f * g) \subset \overline{\supp f + \supp g}\).
}

\exr{}{
Пример того, что:
1)~$*$ не ассоциативна;
2)~$*$ не непрерывна.
}

\rem{Контекст}{
Если \(f, g \in \DD'(\R^n)\) и \(\supp g\) --- компакт, то свёртка
\(f * g\) существует.
}

\exr{}{
Если \(\supp g\) компакт, то \(f * g\) непрерывна по \(f\).
Если \(g_k \to g\) и \(\supp g_k \subset K\) --- компакт, то \(f * g_k \to f * g\).
}

\rem{Контекст}{
\(C \subset \R^n\) --- конус (\(\forall t > 0\ \ tC \subset C\)), причём
\begin{enumerate}[leftmargin=1.6em,itemsep=3pt]
    \item \(C\) выпуклый;
    \item \(C\) замкнутый;
    \item \(C\) открытый, \(\operatorname{Int} C^* \ne \varnothing\),
\end{enumerate}
где \(C^* = \{x \in \R^n \mid (x,y) \ge 0\ \forall y \in C\}\).
Для выпуклых конусов условие~3 равносильно тому, что \(C\) не содержит прямых.

Если \(\supp f \subset C\) и \(\supp g \subset C\), то свёртка \(f * g\)
существует и \(\supp(f * g) \subset C\).
В доказательстве берут \(\eta \equiv 1\) на \(C + B_\epsilon(0)\)
и \(\eta = 0\) вне \(C + B_{2\epsilon}(0)\), и проверяют финитность
\(\eta(x)\eta(y)\, \phi(x+y)\).
От противного найдётся \((x_k, y_k) \to \infty\) с
\[
\eta(x_k)\, \eta(y_k)\, \phi(x_k + y_k) \ne 0,
\]
так что \(x_k, y_k \in C + B_{2\epsilon}\), \(|x_k + y_k| \le A\),
\(\dfrac{|x_k|}{|y_k|} \to 1\) и
\[
\frac{x_k}{|x_k|} \to \hat x \in S_1,
\qquad
\frac{y_k}{|y_k|} \to \hat y \in S_1.
\]
Из \(x_k \in C + B_{2\epsilon}\) следует \(\hat x \in C\).
}

\exr{}{\(\hat y \in C\).}

% =====================================================================
\section{Практика 5}
\renewcommand{\notesheader}{Математическая физика \(\cdot\) Практика 5 \(\cdot\) Упражнения}
% =====================================================================

\rem{Контекст}{
Включение
\(\operatorname{singsupp}(f * g)
\subset
\operatorname{singsupp} f + \operatorname{singsupp} g\)
неверно для свёртки в смысле Владимирова: \(f * g\) существует, если
существует
\[
\lim \langle f(x)\, g(y),\, \eta_k(x,y)\, \phi(x+y) \rangle
\]
и предел не зависит от финитизатора \(\{\eta_k\}\).

По Хёрмандеру \(f * g\) существует, если для всякого компакта
\(K \subset \R^n\) множество
\[
S_K
=
\bigl\{
    (x,y) \in \supp f \times \supp g
    \bigm|
    x + y \in K
\bigr\}
\]
ограничено (значит, компактно).
}

\exr{}{
\(\operatorname{singsupp}(f * g)
\subset
\overline{\operatorname{singsupp} f + \operatorname{singsupp} g}\)
для Хёрмандера.
}

\rem{Контекст}{
Для дифференциального оператора \(L = P(D)\) с постоянными
коэффициентами фундаментальное решение \(\mathcal{E}\) --- это решение
\(L\mathcal{E} = \delta\). Тогда \(u_f = \mathcal{E} * f\) решает \(Lu = f\),
если свёртка существует.

Для обыкновенного оператора
\(\sum_{k=0}^{n} C_k u^{(k)} = f\) фундаментальное решение удовлетворяет
\(\sum_{k=0}^{n} C_k u_{\text{фунд.}}^{(k)} = \delta\),
а общее решение равно \(u_{\text{фунд.}} * f\).
С этой конструкцией связана следующая задача Коши.
}

\exr{}{
\[
\begin{cases}
\sum C_k u^{(k)} = 0, \\
u^{(k)} = 0 \quad \forall k < n-1, \\
u^{(n-1)} = 1.
\end{cases}
\]
}

\rem{Контекст}{
Исходная задача Коши:
\[
\begin{cases}
u_{tt} - a^2 \Delta u = f(x,t), & t > 0, \\
u\big|_{t=0} = \phi, \\
u_t\big|_{t=0} = \psi,
\end{cases}
\qquad
u \in C^2\bigl(\R^n \times \R_{\ge 0}\bigr).
\]
Продолжение \(U\) равно \(u\) при \(t \ge 0\) и нулю при \(t < 0\);
аналогично \(F\) равно \(f\) при \(t \ge 0\) и нулю при \(t < 0\).

Верхнее (обобщённое) уравнение:
\[
U_{tt} - a^2 \Delta U
=
F(x,t) + \phi(x)\, \delta'(t) + \psi(x)\, \delta(t).
\]
}

\exr{}{
Если \(U\) --- решение верхнего уравнения и
\(U \in C^2\bigl(\R^n \times \R_{\ge 0}\bigr)\),
тогда \(u = U\big|_{\{t \ge 0\}}\) --- решение исходного.
}

\end{document}
